\color{unidarkgreen}{\subsection{Максимальная токовая защита стороны СН (МТЗ СН)}}
\color{black}

\begin{enumerate}[label=\arabic{section}.\arabic{subsection}.\arabic*,  align=left, labelsep=4pt, leftmargin=0pt, itemindent=57pt]

\item
Направленная максимальная токовая защита на стороне СН с комбинированным пуском по напряжению
предназначена для резервирования основных защит трансформатора, резервирования отключения КЗ на
шинах СН и на элементах присоединенных к шинам СН.

\begin{figure}[H]
\centering
\includegraphics[width=1\textwidth,height=1\textheight,keepaspectratio]{\fbpath/022.0321 МТЗ СН/img/MTZ_SN.pdf}
\captionof{figure}{Функциональная схема МТЗ СН}\label{pict:MTZ_SN}
\end{figure}

\item
Защита подключается к ТТ, установленным со стороны обмоток СН силового трансформатора.

\item
В состав МТЗ СН входят две направленные ступени с контролем токов трех фаз средней стороны
трансформатора. 

\item
Порог срабатывания первой ступени защиты задается уставкой \textbf{«Iср МТЗ-1 СН»}, а выдержка времени
срабатывания первой ступени регулируется уставкой \textbf{«Тср МТЗ-1 СН»}. Для повышения чувствительности,
предусмотрена возможность пуска первой ступени МТЗ СН от комбинированного пускового органа напряжения
(КПОН) средней стороны трансформатора, которая задается программным переключателем \textbf{«Режим
КПОН МТЗ-1 СН»}. Ввод в работу первой ступени МТЗ СН осуществляется с помощью программного
переключателя \textbf{«МТЗ-1 СН»}. 

\item
Порог срабатывания второй ступени защиты задается уставкой \textbf{«Iср МТЗ-2 СН»}, а выдержка времени
срабатывания второй ступени регулируется уставкой \textbf{«Тср МТЗ-2 СН»}. Для повышения чувствительности,
предусмотрена возможность пуска второй ступени МТЗ СН от комбинированного пускового органа напряжения
(КПОН) средней стороны трансформатора, которая задается программным переключателем \textbf{«Режим
КПОН МТЗ-2 СН»}. Ввод в работу второй ступени МТЗ СН осуществляется с помощью программного
переключателя \textbf{«МТЗ-2 СН»}. 

\item
Оперативный ввод/вывод МТЗ СН производится с помощью БУ (блока управления функциями) \textit{«Вывод
МТЗ СН»}. 

\item
С помощью уставки «ИО МТЗ СН» имеется возможность выбрать по фазным или линейным токам (по
разности токов) будут работать измерительные органы первой и второй ступени МТЗ СН. 

\item
Для обеспечения несрабатывания МТЗ СН от бросков намагничивающего тока предусмотрена
блокировка ступеней МТЗ по содержанию второй гармоники в фазном или линейном токе в зависимости от
уставки \textbf{«ИО МТЗ СН»}. Ввод в работу БНТ МТЗ СН осуществляется с помощью уставки \textbf{«БНТ МТЗ СН»}.
Блокировка основана на контроле отношения действующего значения второй гармоники тока к основной
гармонике. Ступени МТЗ блокируется при превышении данного соотношения пороговой величины, заданной
уставкой \textbf{«3I02г/3I01г МТЗ СН»}. Наличие функции блокировки от бросков намагничивающего тока задается
для каждой ступени отдельно с помощью программных переключателей \textbf{«Режим БНТ МТЗ-1 СН»} и \textbf{«Режим
БНТ МТЗ-2 СН»}. Для предотвращения излишней блокировки при больших токах КЗ, сопровождающихся
насыщением ТТ, блокировка выводится из действия при превышении тока основной гармоники выше уставки
\textbf{«Iср БНТ МТЗ СН»}.

\item
Направленность МТЗ СН обеспечивается органами направления мощности выполненными по 90
-градусной схеме с парами токов и напряжений Ia и Ubc, Ic и Uab. Направление мощности определяется по
величине фазового угла между током и напряжением отдельно для каждой пары сигналов. Разрешение работы
от органа направленности МТЗ СН будет происходить от пары сигналов с максимальным током. За
положительное направление векторов принято вращение против часовой стрелки. Углы максимальной
чувствительности ОНМ откладываются от линейных напряжений и задаются уставками \textbf{«Фмч ОНМ в Т»} и \textbf{«Фмч
ОНМ в сеть»}. Уставки по углам максимальной чувствительности рекомендуется задавать для ОНМ в Т – 45º,
а для ОНМ в сеть – 225º. Ширина зоны срабатывания ±85º от угла максимальной чувствительности. Величина
тока работы ОНМ 0,05$ \cdot I_{\text{ном}}$, величина напряжения работы ОНМ 0,01$ \cdot U_{\text{ном}}$.

\begin{figure}[H]
\centering
\begin{subfigure}{0.47\textwidth}
    \centering
    \includegraphics[width=\textwidth, trim=25mm 13mm 140mm 2mm, clip]{\fbpath/022.0321 МТЗ СН/img/ONM_MTZ_SN.pdf}
    \caption{а)}
    \label{pict_mtz:ONM_MTZ_SN_a}
\end{subfigure}
%\hfill
\begin{subfigure}{0.4\textwidth}
    \centering
    \includegraphics[width=\textwidth, trim=140mm 13mm 35mm 2mm, clip]{\fbpath/022.0321 МТЗ СН/img/ONM_MTZ_SN.pdf}
    \caption{б)}
    \label{pict_mtz:ONM_MTZ_SN_b}
\end{subfigure}

\caption{Функциональная схема ОНМ МТЗ СН}
\label{pict_mtz:ONM_MTZ_SN}
\end{figure}

\item
Для выбора режима направленности по мощности предусмотрены программные переключатели
«Направ МТЗ-1 СН» и «Направ МТЗ-2 СН», имеющий три положения:
\begin{itemize}
\item \textit{«Ненапр»} -- ступень выполнена ненаправленной;
\item \textit{«В Т»} -- ступень выполнена прямонаправленной в сторону Т;
\item \textit{«В сеть»} -- ступень выполнена обратнонаправленной в сторону сети.
\end{itemize}

\item
Имеется возможность задать различные режимы работы ступеней МТЗ СН при возникновении
неисправности в цепях напряжения. Вариант действия логики направленных ступеней МТЗ при срабатывании
КЦН СН задается программными переключателями \textbf{«КЦН ОНМ МТЗ-1 СН»} и \textbf{«КЦН ОНМ МТЗ-2 СН»},
имеющими три положения:
\begin{itemize}
\item \textit{«Выведено»} -- режим работы ступени не зависит от КЦН;
\item \textit{«РНМ»} -- ступень переводится в ненаправленный режим работы;
\item \textit{«МТЗ»} -- при срабатывании КЦН ступень выводится из работы. Однако, если ступень задана ненаправленной, то она останется в работе независимо от срабатывания КЦН. 
\end{itemize}

\item
Предусмотрена возможность оперативного ускорения МТЗ СН для ускорения защит от КЗ на шинах СН
при выводе основной защиты шин. Срабатывание оперативного ускорения МТЗ СН действует на отключение
трансформатора со всех сторон с пуском УРОВ и запретом АПВ. Ввод в работу ОУ МТЗ СН и выбор ускоряемой
ступени осуществляется с помощью программного переключателя \textbf{«ОУ МТЗ СН»}, а выдержка времени
срабатывания ОУ МТЗ СН регулируется уставкой \textbf{«Тср ОУ МТЗ СН»}. Оперативный ввод/вывод ОУ МТЗ СН
производится с помощью команды управления «Ввод ОУ МТЗ СН». 

\item
Направленные ступени МТЗ СН в сторону сети выполнены действующими:
\begin{itemize}
\item c первой выдержкой времени на разделение секции или систем шин СН. Пуск МТЗ СН действует в логику деления средней стороны трансформатора, выдержки времени задаются уставками \textbf{«Тср МТЗ-1 откл СВ(ШСВ)»} и \textbf{«Тср МТЗ-2 откл СВ(ШСВ)»};
\item cо второй выдержкой времени на отключение (с пуском УРОВ и запретом АПВ) выключателей стороны СН. Пуск МТЗ СН действует в логику деления средней стороны трансформатора, выдержки времени задаются уставками \textbf{«Тср МТЗ-1 откл В»} и \textbf{«Тср МТЗ-2 откл В»};
\item c выдержкой времени срабатывание МТЗ СН на отключение трансформатора со всех сторон с пуском УРОВ и запретом АПВ. Выдержки времени задаются уставками \textbf{«Тср МТЗ-1 СН»} и \textbf{«Тср МТЗ-2 СН»}.
\end{itemize}

\item
Действие срабатывания первой и второй ступени МТЗ СН предусматривает отключение
трансформатора со всех сторон с пуском УРОВ и запретом АПВ. 

\small
\setlength{\extrarowheight}{0.05cm} % высота ячеек в таблице
\begin{longtable}
{|>{\centering\arraybackslash}m{4.5cm}|>{\centering\arraybackslash}m{4.5cm}|>{\centering\arraybackslash}m{3cm}|>{\centering\arraybackslash}m{3.5cm}|}
\caption{Перечень уставочных параметров МТЗ СН\hfill\vspace{-0.5\baselineskip}}\label{mtz_sn:tbl1} \\ 
\hhline{|-|-|-|-|}
\rowcolor{gray!30}
Наименование уставки & Диапазон & Шаг &  Значение по умолчанию \\ 
\hhline{|-|-|-|-|}
\endfirsthead

\caption*{\hspace{3pt}\emph{Продолжение таблицы \ref{mtz_sn:tbl1}\hfill\vspace{-0.5\baselineskip}}} \\ % сделано по ГОСТ 2.105 п.6.8.7
\hhline{|-|-|-|-|}
\rowcolor{gray!30}
Наименование уставки & Диапазон & Шаг &  Значение по умолчанию \\  
\hhline{|-|-|-|-|}
\endhead
\hhline{|-|-|-|-|}

\endfoot
\endlastfoot

%%%%%%%%%%%%содержимое%%%%%%%%%%%%%%

\end{longtable}
\setlength{\extrarowheight}{0.0cm} % возврат параметра высота ячеек в таблице
\normalsize










































\end{enumerate} 